\scanpage{100}
\begin{center}- 12 -\end{center}
\margin{n° 285}

se borner au cas où $I$ est un intervalle de \qquad, et comme la proposition est triviale si $I$ est fini, on peut supposer $I=\unsure{\mathbf N|\mathbb N}$. Soit $x_p\in G_p$; l'hypothèse permet de définir par récurrence sur $n\geq p$ une suite $(x_n)_{n\geq p}$ satisfaisant aux conditions $x_n\in G_n$, $f_{n,n+1}(x_{n+1})=x_n$; il est clair alors que si on pose, pour $m<p$, $x_m=f_{mp}(x_p)$ la suite $(x_n)_{n\geq0}=x$ appartient à $G$ et que $f_p(x)=x_p$, d'où la proposition.

PROPOSITION 23 - \uline{Les notations étant celles du cor.1 de la prop.20, on suppose en outre que pour tout} $\alpha\in I$, \uline{l'application canonique} $\varphi_\alpha$ \uline{de} $A$ \uline{dans} $A_\alpha$ \uline{soit surjective}. \uline{Alors, si pour tout} $\alpha\in I$ $g_\alpha(A_\alpha)$ \uline{est dense dans} $B_\alpha$, $g(A)$ \uline{est dense dans} $B$.

En effet, il suffit de montrer que, pour tout $\alpha\in I$ et tout ensemble ouvert non vide $U_\alpha$ dans $B_\alpha$, il existe $x\in A$ tel que $\psi_\alpha(g(x))\in U_\alpha$ (en désignant par $\psi_\alpha$ l'application canonique de $B$ dans $B_\alpha$). Mais par hypothèse il existe $x_\alpha\in A_\alpha$ tel que $g_\alpha(x_\alpha)\in U_\alpha$, et il existe $x\in A$ tel que $\varphi_\alpha(x)=x_\alpha$; comme $\psi_\alpha(g(x))=g_\alpha(\varphi_\alpha(x))$, cet élément répond à la question

PROPOSITION 24 - \uline{Soient} $(G_\alpha,f_{\alpha\beta})$ \uline{un système projectif de groupes topologiques}, $G=\varprojlim G_\alpha$, $f_\alpha$ \uline{l'application canonique de} $G$ \uline{dans} $G_\alpha$. \uline{On suppose que} : 1°) \uline{pour tout} $\alpha\in I$, $f_\alpha$ \uline{est surjective} ; 2°) \uline{pour tout couple d'indices tel que} $\alpha\leq\beta$, $f_{\alpha\beta}$ \uline{est un homomorphisme de} $G_\beta$ \uline{sur} $G_\alpha$. \uline{Dans ces conditions}, $f_\alpha$ \uline{est un homomorphisme de} $G$ \uline{sur} $G_\alpha$.

En effet, soit $V$ un voisinage quelconque de $e$ dans $G$. Il existe un indice $\beta\geq\alpha$ et un voisinage $V_\beta$ de $e_\beta$ dans $G_\beta$ tel que $W=f_\beta^{-1}(V_\beta)$ soit un voisinage de $e$ contenu dans $V$. Comme $f_\beta$ est surjective, on a $V_\beta=f_\beta(W)$, et comme $f_{\alpha\beta}$ est un homomorphisme surjectif, $f_{\alpha\beta}(V_\beta)$ est un voisinage de $e_{\unsure{\alpha}}$ dans $G_\alpha$; mais on a $f_\alpha(W)=f_{\alpha\beta}(f_\beta(W))=f_{\alpha\beta}(V_\beta)$ d'où la proposition (th.1).

Soit $N_\alpha$ le noyau de $f_\alpha$; les $N_\alpha$ constituent une base de filtre, formée

\margin{99}
